\section{Direct Products and Finitely Generated Abelian Groups}
\textbf{Conventions.} Permutations compose from right to left. The identity is $e$. A group generated by an element $a$ is denoted $\gen{a}$.

\exercise{11.1}{Elements and orders in $\Z_{2}\times\Z_{4}$}
The elements are all pairs $(a,b)$ with $0\le a<2$ and $0\le b<4$. Addition is componentwise. A multiple $k(a,b)$ is zero exactly when both $ka\equiv0\pmod{2}$ and $kb\equiv0\pmod{4}$. Hence its order is the least common multiple of the two component orders. The table lists the order of each pair, with first coordinate indexing rows and second coordinate indexing columns:\[\begin{array}{c|cccc}a\backslash b&0&1&2&3\\\hline0 & 1 & 4 & 2 & 4\\
1 & 2 & 4 & 2 & 4\\\end{array}\]No element has order $8$, so the order-$8$ group is not cyclic.

\exercise{11.2}{Elements and orders in $\Z_{3}\times\Z_{4}$}
The elements are all pairs $(a,b)$ with $0\le a<3$ and $0\le b<4$. Addition is componentwise. A multiple $k(a,b)$ is zero exactly when both $ka\equiv0\pmod{3}$ and $kb\equiv0\pmod{4}$. Hence its order is the least common multiple of the two component orders. The table lists the order of each pair, with first coordinate indexing rows and second coordinate indexing columns:\[\begin{array}{c|cccc}a\backslash b&0&1&2&3\\\hline0 & 1 & 4 & 2 & 4\\
1 & 3 & 12 & 6 & 12\\
2 & 3 & 12 & 6 & 12\\\end{array}\]The element $(1,1)$ has order $\lcm(3,4)=12$, equal to the group order, so it generates the whole group; the group is cyclic.

\exercise{11.3}{Order of $(2,6)$}
In $\Z_m$, the order of $a$ is $m/\gcd(m,a)$. In the given group $\Z_{4}\times\Z_{12}$, the component orders are\[\frac{4}{\gcd(4,2)}=2,\quad \frac{12}{\gcd(12,6)}=2.\]All components return to zero simultaneously at the least common multiple of these orders, so\[\operatorname{ord}(2,6)=\lcm(2,2)=2.\]

\exercise{11.4}{Order of $(2,3)$}
In $\Z_m$, the order of $a$ is $m/\gcd(m,a)$. In the given group $\Z_{6}\times\Z_{15}$, the component orders are\[\frac{6}{\gcd(6,2)}=3,\quad \frac{15}{\gcd(15,3)}=5.\]All components return to zero simultaneously at the least common multiple of these orders, so\[\operatorname{ord}(2,3)=\lcm(3,5)=15.\]

\exercise{11.5}{Order of $(8,10)$}
In $\Z_m$, the order of $a$ is $m/\gcd(m,a)$. In the given group $\Z_{12}\times\Z_{18}$, the component orders are\[\frac{12}{\gcd(12,8)}=3,\quad \frac{18}{\gcd(18,10)}=9.\]All components return to zero simultaneously at the least common multiple of these orders, so\[\operatorname{ord}(8,10)=\lcm(3,9)=9.\]

\exercise{11.6}{Order of $(3,10,9)$}
In $\Z_m$, the order of $a$ is $m/\gcd(m,a)$. In the given group $\Z_{4}\times\Z_{12}\times\Z_{15}$, the component orders are\[\frac{4}{\gcd(4,3)}=4,\quad \frac{12}{\gcd(12,10)}=6,\quad \frac{15}{\gcd(15,9)}=5.\]All components return to zero simultaneously at the least common multiple of these orders, so\[\operatorname{ord}(3,10,9)=\lcm(4,6,5)=60.\]

\exercise{11.7}{Order of $(3,6,12,16)$}
In $\Z_m$, the order of $a$ is $m/\gcd(m,a)$. In the given group $\Z_{4}\times\Z_{12}\times\Z_{20}\times\Z_{24}$, the component orders are\[\frac{4}{\gcd(4,3)}=4,\quad \frac{12}{\gcd(12,6)}=2,\quad \frac{20}{\gcd(20,12)}=5,\quad \frac{24}{\gcd(24,16)}=3.\]All components return to zero simultaneously at the least common multiple of these orders, so\[\operatorname{ord}(3,6,12,16)=\lcm(4,2,5,3)=60.\]

\exercise{11.8}{Largest orders of cyclic subgroups}
In $\Z_m\times\Z_n$, every element order divides $\lcm(m,n)$, since that many additions kill both coordinates. The element $(1,1)$ attains this order, so it is the largest cyclic-subgroup order. Thus the answers are
\[\lcm(6,8)=24,\qquad\lcm(12,15)=60.\]
These need not equal the group orders $48$ and $180$, respectively.

\exercise{11.9}{Proper nontrivial subgroups of $\Z_2^2$}
All three nonzero elements have order $2$. They generate the three subgroups
\[\{(0,0),(1,0)\},\quad\{(0,0),(0,1)\},\quad\{(0,0),(1,1)\}.\]
Any subgroup containing two different nonzero vectors contains their sum, the third nonzero vector, and so is the whole group. Thus the three listed subgroups are exactly the proper nontrivial ones.

\exercise{11.10}{Proper nontrivial subgroups of $\Z_2^3$}
Regard the group as the three-dimensional vector space over $\Z_2$; subgroup closure is exactly linear-subspace closure, because the only scalars are $0,1$. The seven subgroups of order $2$ are $\{0,v\}$ for
\[v=(1,0,0),(0,1,0),(0,0,1),(1,1,0),(1,0,1),(0,1,1),(1,1,1).\]
The seven subgroups of order $4$ are the planes described, respectively, by
\[x=0,\quad y=0,\quad z=0,\quad x+y=0,\quad x+z=0,\quad y+z=0,\quad x+y+z=0,\]
where arithmetic is modulo $2$. Each displayed equation has four solutions. Every two-dimensional subspace is the kernel of one nonzero linear functional, and over $\Z_2$ no two distinct nonzero functionals are scalar multiples. Consequently this lists all seven planes. Proper nontrivial subgroups have dimension $1$ or $2$, so these fourteen are all of them.

\exercise{11.11}{Order-four subgroups of $\Z_2	imes\Z_4$}
An element has order $4$ precisely when its second coordinate is odd. There are four such elements, and each cyclic order-$4$ subgroup has two generators. The two distinct cyclic subgroups are
\[\gen{(0,1)}=\{(0,0),(0,1),(0,2),(0,3)\},\]
\[\gen{(1,1)}=\{(0,0),(1,1),(0,2),(1,3)\}.\]
Any noncyclic group of order $4$ is a Klein group and all its elements are killed by $2$. The elements killed by $2$ here are exactly
\[\{(0,0),(1,0),(0,2),(1,2)\},\]
which form one more subgroup. This exhausts both possible types of order-$4$ subgroup.

\exercise{11.12}{Klein subgroups of $\Z_2	imes\Z_4	imes\Z_4$}
Every element of a Klein subgroup is killed by $2$, so it must lie in
\[T=\{(a,2b,2c):a,b,c\in\Z_2\}\cong\Z_2^3.\]
A Klein subgroup is a two-dimensional subspace of this three-dimensional space. Thus all seven are given by $(a,2b,2c)$ subject to one of the conditions
\[a=0,\ b=0,\ c=0,\ a+b=0,\ a+c=0,\ b+c=0,\ a+b+c=0\]
in $\Z_2$. For example the first is $\{(0,0,0),(0,2,0),(0,0,2),(0,2,2)\}$. Each condition gives four elements, all of exponent $2$, hence a Klein group. The seven-plane classification from the preceding exercise proves completeness.

\exercise{11.13}{Direct product descriptions of $\Z_{60}$}
A finite direct product of cyclic groups is cyclic precisely when their orders are pairwise relatively prime. Since $60=4\cdot3\cdot5$, the prime-power blocks $4,3,5$ must remain intact but can be grouped. Ignoring order of factors, all descriptions using at least two nontrivial factors are
\[\Z_4\times\Z_{15},\quad\Z_3\times\Z_{20},\quad\Z_5\times\Z_{12},\quad\Z_4\times\Z_3\times\Z_5.\]
These exhaust the partitions of the three prime-power blocks into at least two groups. If factors $\Z_1$ are permitted, one can insert arbitrarily many trivial factors, giving infinitely many formally different descriptions without changing the group; the finite list uses the usual convention of omitting trivial factors.

\exercise{11.14}{Fill in the numerical blanks}
\begin{enumerate}[label=(\alph*)]
\item $\operatorname{ord}_{\Z_{24}}(18)=24/\gcd(24,18)=24/6=4$.
\item $|\Z_3\times\Z_4|=3\cdot4=12$.
\item The orders of $4\in\Z_{12}$ and $2\in\Z_8$ are $3$ and $4$, so $(4,2)$ has order $12$.
\item $V\cong\Z_2\times\Z_2$; both blanks are $2$.
\item A finite-order element of $\Z_2\times\Z\times\Z_4$ must have zero middle coordinate. There are $2\cdot1\cdot4=8$ such elements, all indeed of finite order.
\end{enumerate}

\exercise{11.15}{Maximum order in $\Z_4	imes\Z_6$}
Every component is killed by its modulus, so $12=\lcm(4,6)$ kills every pair. The pair $(1,1)$ has component orders $4,6$ and hence order $12$. Thus $12$ is the maximum possible order.

\exercise{11.16}{Comparing two groups of order $24$}
Yes. Splitting relatively prime cyclic factors gives
\[\Z_2\times\Z_{12}\cong\Z_2\times\Z_4\times\Z_3,
\qquad\Z_4\times\Z_6\cong\Z_4\times\Z_2\times\Z_3.\]
The right-hand products differ only in order of coordinates, so they are isomorphic. An explicit map from the second group to the first is $(a,b)\mapsto(b\bmod2,c)$, where $c$ is the unique residue modulo $12$ satisfying $c\equiv a\pmod4$, $c\equiv b\pmod3$; existence and uniqueness follow from coprimality of $4,3$.

\exercise{11.17}{Maximum order in a three-factor product}
For $\Z_8\times\Z_{10}\times\Z_{24}$ the maximum is $\lcm(8,10,24)=2^3\cdot3\cdot5=120$. Every element is killed by $120$, and $(1,1,1)$ has exactly that order, proving both the upper bound and its attainment.

\exercise{11.18}{Comparing the groups in Exercise 18}
The groups are not isomorphic, although both have the same order and the same maximum element order $120$. Count elements killed by $4$. In $\Z_m$ the congruence $4x=0$ has $\gcd(4,m)$ solutions. Therefore the counts are
\[\gcd(4,8)\gcd(4,10)\gcd(4,24)=4\cdot2\cdot4=32,\]
\[\gcd(4,4)\gcd(4,12)\gcd(4,40)=4\cdot4\cdot4=64.\]
An isomorphism preserves the equation $4x=0$ and hence its number of solutions, so the different counts rule out an isomorphism.

\exercise{11.19}{Maximum order in $\Z_4	imes\Z_{18}	imes\Z_{15}$}
The least common multiple of the factor orders is $\lcm(4,18,15)=2^2\cdot3^2\cdot5=180$. It kills every element, and $(1,1,1)$ has that order. Therefore the maximum is $180$.

\exercise{11.20}{Comparing the groups in Exercise 20}
Yes. Decompose the factors into prime powers:
\[\begin{aligned}
\Z_4\times\Z_{18}\times\Z_{15}&\cong\Z_4\times\Z_2\times\Z_9\times\Z_3\times\Z_5,\\
\Z_3\times\Z_{36}\times\Z_{10}&\cong\Z_3\times\Z_4\times\Z_9\times\Z_2\times\Z_5.
\end{aligned}\]
The same five prime-power cyclic factors occur, so reordering coordinates gives an isomorphism. Splitting was legitimate because each original factor was split into relatively prime factors.

\exercise{11.21}{Abelian groups of order $8$}
The classification theorem assigns one group to each partition of $3$, using one cyclic factor $\Z_{2^k}$ for each part $k$. Thus the complete list is\[\begin{gathered}\Z_{8},\\\Z_{4}\times\Z_{2},\\\Z_{2}\times\Z_{2}\times\Z_{2}.\end{gathered}\]The partitions are distinct, so the uniqueness in the classification theorem shows the listed groups are pairwise nonisomorphic. There are $3$ groups in all.

\exercise{11.22}{Abelian groups of order $16$}
The classification theorem assigns one group to each partition of $4$, using one cyclic factor $\Z_{2^k}$ for each part $k$. Thus the complete list is\[\begin{gathered}\Z_{16},\\\Z_{8}\times\Z_{2},\\\Z_{4}\times\Z_{4},\\\Z_{4}\times\Z_{2}\times\Z_{2},\\\Z_{2}\times\Z_{2}\times\Z_{2}\times\Z_{2}.\end{gathered}\]The partitions are distinct, so the uniqueness in the classification theorem shows the listed groups are pairwise nonisomorphic. There are $5$ groups in all.

\exercise{11.23}{Abelian groups of order $32$}
The classification theorem assigns one group to each partition of $5$, using one cyclic factor $\Z_{2^k}$ for each part $k$. Thus the complete list is\[\begin{gathered}\Z_{32},\\\Z_{16}\times\Z_{2},\\\Z_{8}\times\Z_{4},\\\Z_{8}\times\Z_{2}\times\Z_{2},\\\Z_{4}\times\Z_{4}\times\Z_{2},\\\Z_{4}\times\Z_{2}\times\Z_{2}\times\Z_{2},\\\Z_{2}\times\Z_{2}\times\Z_{2}\times\Z_{2}\times\Z_{2}.\end{gathered}\]The partitions are distinct, so the uniqueness in the classification theorem shows the listed groups are pairwise nonisomorphic. There are $7$ groups in all.

\exercise{11.24}{Abelian groups of order $720$}
Factor $720=2^4\cdot3^2\cdot5$. The possible $2$-primary factors are
\[T\in\{\Z_{16},\ \Z_8\times\Z_2,\ \Z_4\times\Z_4,\ \Z_4\times\Z_2^2,\ \Z_2^4\},\]
and the possible $3$-primary factors are $U\in\{\Z_9,\Z_3^2\}$. The $5$-primary factor must be $\Z_5$. All groups, with none omitted or repeated up to isomorphism, are therefore
\[\boxed{T\times U\times\Z_5\quad\text{for every one of the five }T\text{ and two }U.}\]
This specifies $5\cdot2=10$ groups explicitly by independent choices; uniqueness of the prime-power cyclic decomposition distinguishes them.

\exercise{11.25}{Abelian groups of order $1089$}
Since $1089=33^2=3^2\cdot11^2$, each primary component has two possibilities. The four groups are
\[\begin{gathered}
\Z_9\times\Z_{121},\qquad\Z_9\times\Z_{11}^2,\\
\Z_3^2\times\Z_{121},\qquad\Z_3^2\times\Z_{11}^2.
\end{gathered}\]
The exponent-$2$ partition choices are $2$ and $1+1$ for each prime. They produce all primary decompositions and hence all isomorphism types.

\exercise{11.26}{Counts for orders $24$, $25$, and $600$}
The printed numbers are $24$ and $25$, not $2^4$ and $2^5$. Since $24=2^3\cdot3$, the three partitions of $3$ give three groups:
\[\Z_8\times\Z_3,\quad\Z_4\times\Z_2\times\Z_3,\quad\Z_2^3\times\Z_3.\]
For order $25=5^2$ the two possibilities are $\Z_{25}$ and $\Z_5^2$. Because $24$ and $25$ are coprime, their primary decompositions can be chosen independently. Hence for $(24)(25)=600$ there are $3\cdot2=6$ abelian groups up to isomorphism.

\exercise{11.27}{Multiplicativity for coprime group orders}
Suppose $\gcd(m,n)=1$. The prime factors of $m$ and $n$ are disjoint. By the finite abelian classification theorem, any abelian group of order $mn$ is uniquely determined by its primary factors. Collecting the factors at primes dividing $m$ gives a group of order $m$; collecting those at primes dividing $n$ gives one of order $n$. Conversely, the direct product of any one of the $r$ types of order $m$ and any one of the $s$ types of order $n$ gives a group of order $mn$. Uniqueness of primary decompositions means different pairs give different isomorphism types. Thus there are exactly $rs$.

\exercise{11.28}{Groups of order $10^5$}
Factor $10^5=2^5\cdot5^5$. There are seven partitions of $5$, hence seven abelian $2$-groups of order $2^5$ and seven abelian $5$-groups of order $5^5$. These choices are independent by coprimality, so the answer is $7\cdot7=49$.

\exercise{11.29}{Partition counts for prime-power orders}
\textbf{(a)} An abelian group of order $p^n$ corresponds to a partition of $n$: parts $\lambda_1,\ldots,\lambda_k$ give $\Z_{p^{\lambda_1}}\times\cdots\times\Z_{p^{\lambda_k}}$. The required table is
\[\begin{array}{c|rrrrrrr}n&2&3&4&5&6&7&8\\\hline\text{number}&2&3&5&7&11&15&22\end{array}\]
One reproducible way to obtain it is the recurrence $P(n,k)=P(n,k-1)+P(n-k,k)$, splitting partitions according to whether a part $k$ occurs, with $P(0,k)=1$ and $P(n,0)=0$ for $n>0$; the desired count is $P(n,n)$.
\textbf{(b)} Distinct primes permit independent choices. Thus
\[\begin{aligned}
\text{(i) }p^3q^4r^7 &: 3\cdot5\cdot15=225,\\
\text{(ii) }(qr)^7=q^7r^7 &:15\cdot15=225,\\
\text{(iii) }q^5r^4q^3=q^8r^4 &:22\cdot5=110.
\end{aligned}\]
In (iii) the two powers of the same prime $q$ must first be combined; they are not independent primary factors.

\exercise{11.30}{A Cayley digraph on a toroidal grid}
Place a vertex $(i,j)$ in an $m$-by-$n$ grid, with $0\le i<m$, $0\le j<n$. Draw one solid arrow
\[(i,j)\longrightarrow(i+1\bmod m,j)\]
for the generator $(1,0)$, and one dashed arrow
\[(i,j)\dashrightarrow(i,j+1\bmod n)\]
for $(0,1)$. Thus each row wraps into an $m$-cycle and each column into an $n$-cycle. A corner of the schematic is
\[\begin{array}{ccccc}
(0,2)&\longrightarrow&(1,2)&\longrightarrow&\cdots\\
\uparrow&&\uparrow&&\\
(0,1)&\longrightarrow&(1,1)&\longrightarrow&\cdots\\
\uparrow&&\uparrow&&\\
(0,0)&\longrightarrow&(1,0)&\longrightarrow&\cdots
\end{array}\]
Vertical arrows represent the second arc type; include the wrap arrows from $i=m-1$ to $i=0$ and from $j=n-1$ to $j=0$. These rules specify every labeled vertex and arc, including the boundary identifications.

\exercise{11.31}{Two concentric polygon digraphs}
Let $r$ denote the solid-arrow generator, of order $n$, and $s$ the dashed-edge generator, of order $2$. Label one polygon $e,r,\ldots,r^{n-1}$ and the matching vertices in the other $s,rs,\ldots,r^{n-1}s$.
\textbf{(a)} The group is abelian exactly when the arrows on the two polygons run in the same direction. In that case traversing a dashed edge before or after a solid edge reaches the same vertex, so $rs=sr$; because $r,s$ generate the group, this proves commutativity. Opposite directions instead give $sr=r^{-1}s\ne rs$ for $n\ge3$.
\textbf{(b)} In the abelian case all $2n$ elements are uniquely $r^is^j$ with $0\le i<n$, $j=0,1$, and multiplication adds both exponents. Thus $G\cong\Z_n\times\Z_2$.
\textbf{(c)} This product is cyclic exactly when $\gcd(n,2)=1$, that is, $n$ is odd.
\textbf{(d)} Opposite arrow directions give exactly the usual dihedral relations, and hence the $2n$-element group is $D_n$.

\exercise{11.32}{True or false: direct products}
\begin{enumerate}[label=(\alph*)]
\item True: $(g_1,g_2)\mapsto(g_2,g_1)$ is an isomorphism.
\item Componentwise computation is precisely the multiplication rule; the description ``easy'' is informal, but the intended mathematical content is true.
\item False: $\Z\times\Z$ is a direct product of infinite groups.
\item True: the orders of two proper nontrivial finite factors would both exceed $1$ and multiply to the prime group order, which is impossible.
\item False: every element of $\Z_2\times\Z_4$ has order dividing $4$, unlike a generator of $\Z_8$.
\item False: the product has $8$ elements, while $S_8$ has $8!$; it is also abelian.
\item False: both groups have $24$ elements, but $\Z_3\times\Z_8$ is abelian and $S_4$ is not.
\item False: the identity has order $1$; other orders can also be $2$ or $4$.
\item False: $|\Z_{12}\times\Z_{15}|=12\cdot15=180$.
\item True: there are independently $m$ choices in the first coordinate and $n$ in the second.
\end{enumerate}

\exercise{11.33}{An indecomposable abelian group}
Take $G=\Z_p$ for a prime $p$. Its only subgroups are $\{0\}$ and itself, so there are no proper nontrivial subgroups to use as internal factors. Alternatively, an internal product of two nontrivial finite groups would have composite order, contrary to $|G|=p$.

\exercise{11.34}{Subgroups isomorphic to the full product}
\textbf{(a)} Exactly one. The finite group $\Z_5\times\Z_6$ has $30$ elements, and a subgroup isomorphic to it must also have $30$ elements, hence equal the whole group.
\textbf{(b)} Infinitely many: for each positive integer $k$, $k\Z\times\Z$ is a subgroup of $\Z^2$ isomorphic to $\Z^2$ via $(a,b)\mapsto(ka,b)$. These subgroups are distinct. More precisely there are countably infinitely many, since any subgroup isomorphic to $\Z^2$ has two generators in the countable set $\Z^2$, and there are only countably many ordered pairs of possible generators.

\exercise{11.35}{A composite-order group with no nontrivial direct decomposition}
Take $\Z_4$. Its only proper nontrivial subgroup is $\{0,2\}$. Any two proper nontrivial internal factors would therefore coincide and have nontrivial intersection, so cannot give an internal direct product. Equivalently the only possible nontrivial factor orders are $2,2$, whose external product is $\Z_2^2$, with no element of order $4$, and hence is not $\Z_4$.

\exercise{11.36}{True or false: classification of abelian groups}
\begin{enumerate}[label=(\alph*)]
\item True: every group of prime order is cyclic.
\item False: $\Z_2^2$ has prime-power order $4$ and is not cyclic.
\item False: every sum of $4$'s and $6$'s modulo $8$ is even, so these generate only $\{0,2,4,6\}$.
\item True: $5-4=1$, so $\{4,5,6\}$ generates $1$ and hence $\Z_8$.
\item True: a finite group is finitely generated, so the theorem applies.
\item False: $\Z$ and $\Z\times\Z_2$ both have Betti number $1$, but only the latter has nontrivial torsion.
\item True: the primary decomposition contains a $5$-power cyclic factor, which has an element of order $5$.
\item False: $\Z_2^2$ has order divisible by $4$ but no element of order $4$.
\item True: the primary decomposition contains elements of orders $2$ and $3$ in different primary factors; their product has order $6$.
\item True: a nonzero free $\Z$ factor would make the group infinite, so a finite group's Betti number is $0$.
\end{enumerate}

\exercise{11.37}{Changing the prime in a prime-power order}
The numbers are equal. For either prime, the isomorphism types are parametrized by the same partitions of the exponent $r$. Replacing every factor $\Z_{p^k}$ in a decomposition by $\Z_{q^k}$ gives a bijection of the two lists of isomorphism types. Thus the count depends on $r$ alone, not on the prime.

\exercise{11.38}{Order-eight and order-four subgroups of an order-$72$ group}
Write $G=P\times Q$ with $|P|=8$ and $|Q|=9$, using the primary decomposition.
\textbf{(a)} There is exactly one subgroup of order $8$, namely $P\times\{0\}$. An element of a group of order $8$ has order a power of $2$, so its projection into the order-$9$ group $Q$ is zero. Thus any order-$8$ subgroup lies in $P\times\{0\}$ and must equal it.
\textbf{(b)} The number of order-$4$ subgroups is not determined by $|G|$ alone. If $P=\Z_8$, it has exactly one order-$4$ subgroup. If $P=\Z_2^3$, it has seven planes of order $4$. Taking, for example, $Q=\Z_9$ gives two groups of order $72$ with different answers.

\exercise{11.39}{The torsion subgroup}
Let $T$ consist of elements of finite order. The identity is in $T$. If $a^r=e$ and $b^s=e$, commutativity gives $(ab)^{rs}=a^{rs}b^{rs}=e$, so $ab$ has finite order and belongs to $T$. If $a^r=e$, then $(a^{-1})^r=(a^r)^{-1}=e$, so $a^{-1}\in T$. Thus $T$ is a subgroup. The commutativity assumption is used in expanding the power of $ab$; it cannot simply be omitted.

\exercise{11.40}{Orders of two torsion subgroups}
In a product with a $\Z$ coordinate, finite order forces that coordinate to be $0$, because a nonzero integer has infinite additive order. All elements in the remaining finite factors are torsion. Hence the torsion subgroups are
\[\Z_4\times\{0\}\times\Z_3\quad\text{of order }12,\]
\[\Z_{12}\times\{0\}\times\Z_{12}\quad\text{of order }144.\]

\exercise{11.41}{Torsion in the nonzero real numbers}
If $x^n=1$ for a real nonzero $x$, then $|x|^n=1$, so $|x|=1$ and $x=\pm1$. Conversely $1$ has order $1$ and $-1$ has order $2$. Thus the torsion subgroup is $\{1,-1\}\cong\Z_2$.

\exercise{11.42}{Torsion in the nonzero complex numbers}
A complex number has finite multiplicative order exactly when it is a root of unity. Indeed $z^n=1$ implies $z=e^{2\pi i k/n}$ for some integer $k$, and any such number has a finite power equal to $1$. Therefore
\[T=\{e^{2\pi i q}:q\in\Q\},\]
the group of all roots of unity. It is isomorphic to $\Q/\Z$ under the map $q+\Z\mapsto e^{2\pi iq}$.

\exercise{11.43}{Torsion and free parts of a finitely generated abelian group}
By the classification theorem choose an isomorphism $G\cong F\times\Z^r$, where $F$ is a finite product of finite cyclic prime-power groups. An element has finite order precisely when all its $\Z$ coordinates are zero, so the torsion subgroup corresponds to $F\times\{0\}$. The subgroup corresponding to $\{e\}\times\Z^r$ is torsion free. These subgroups intersect trivially, commute elementwise, and every element is their product; transporting them back through the isomorphism gives an internal direct-product decomposition of $G$. This includes the case of a trivial finite or free part.

\exercise{11.44}{Torsion coefficients}
\textbf{(a)} Since $\gcd(4,9)=1$, $\Z_4\times\Z_9\cong\Z_{36}$. The only torsion coefficient is $36$.
\textbf{(b)} The group is $\Z_6\times\Z_{12}\times\Z_{20}$. Split into prime-power factors and sort the exponents at each prime, padding with zero exponents on the left:
\[\begin{array}{c|ccc}p&\text{first}&\text{second}&\text{third}\\\hline2&1&2&2\\3&0&1&1\\5&0&0&1\end{array}\]
Combining each column gives $2$, $2^2\cdot3=12$, and $2^2\cdot3\cdot5=60$. Hence the torsion coefficients are $2,12,60$, and $2\mid12\mid60$.
\textbf{(c)} For an arbitrary finite product, factor all moduli into prime powers. For each prime put the exponents in nondecreasing order and pad shorter lists with zero exponents at their left ends to obtain equal lengths. Multiply the prime powers column by column. The resulting column products divide the next ones, because every prime exponent is nondecreasing. Reordering primary factors and combining relatively prime factors proves this process preserves the isomorphism type. Discard any leading coefficient $1$.

\exercise{11.45}{Synopsis of the coprime-product criterion}
The element $(1,1)\in\Z_m\times\Z_n$ has order $\lcm(m,n)$, so when $\gcd(m,n)=1$ it generates the entire order-$mn$ group. If the gcd is $d>1$, every element is killed by $mn/d<mn$, so no element can generate the whole group and the product is not cyclic.

\exercise{11.46}{A direct product of abelian groups is abelian}
Let $a=(a_1,\ldots,a_r)$ and $b=(b_1,\ldots,b_r)$ belong to the product. Since each factor is abelian, $a_ib_i=b_ia_i$ for every $i$. Therefore
\[ab=(a_1b_1,\ldots,a_rb_r)=(b_1a_1,\ldots,b_ra_r)=ba.\]
This is exactly commutativity of the product operation.

\exercise{11.47}{Elements of order at most two}
The given set is $H=\{g\in G:g^2=e\}$. It contains $e$. If $a,b\in H$, commutativity gives $(ab)^2=a^2b^2=e$, so $ab\in H$. Also $(a^{-1})^2=(a^2)^{-1}=e$. Thus $H$ is a subgroup. Any nonidentity element killed by $2$ has order exactly $2$, so this is precisely the set described in the question.

\exercise{11.48}{Identity together with all elements of an exact order}
For order $3$, the answer is yes: the set is $\{g:g^3=e\}$, and commutativity shows it is closed under products and inverses. For order $4$, the answer is no. In $\Z_4$, $1$ has order $4$ but $1+1=2$ has order $2$, so the set $\{0,1,3\}$ is not closed.
In general the universal property holds exactly for $n=1$ and for prime $n$. For prime $p$, every nonidentity solution of $g^p=e$ has order $p$, so the same kernel-style closure argument works. If $n$ is composite, choose a prime $p\mid n$ with $p<n$ and work in $\Z_n\times\Z_n$. Both
\[a=(1,0),\qquad b=(-1,n/p)\]
have order $n$ because their first coordinates do. But $a+b=(0,n/p)$ has order $p$, which is neither $1$ nor $n$. Thus the prescribed set fails closure. This also explains the distinction between exact order $n$ and order dividing $n$.

\exercise{11.49}{Why commutativity is needed for involutions}
In the nonabelian group $S_3$, the identity and three transpositions are exactly the elements of order $1$ or $2$. But $(1,2)(2,3)=(1,2,3)$ has order $3$ and does not lie in that set. Hence it is not a subgroup, providing a counterexample when the abelian hypothesis is removed.

\exercise{11.50}{The coordinate subgroups in an external product}
Write the two coordinate subgroups as $\widetilde H=H\times\{e_K\}$ and $\widetilde K=\{e_H\}\times K$.
\textbf{(a)} Every $(h,k)$ equals $(h,e_K)(e_H,k)$.
\textbf{(b)} Componentwise multiplication gives
\[(h,e_K)(e_H,k)=(h,k)=(e_H,k)(h,e_K),\]
so the two subgroups commute elementwise even if the individual groups are nonabelian.
\textbf{(c)} An element in both coordinate subgroups has first coordinate $e_H$ and second coordinate $e_K$, so the intersection is exactly $\{(e_H,e_K)\}$.

\exercise{11.51}{Recognizing an internal direct product}
By the first assumed condition, every $g\in G$ is $hk$. If $hk=h'k'$, rearrange to obtain
\[h'^{-1}h=k'k^{-1}\in H\cap K=\{e\}.
\]
Thus $h=h'$ and $k=k'$, proving uniqueness. Define $\Phi:H\times K\to G$ by $\Phi(h,k)=hk$. Existence and uniqueness say it is bijective. Since every element of $H$ commutes with every element of $K$,
\[\Phi(h,k)\Phi(h',k')=hkh'k'=hh'kk'=\Phi(hh',kk').\]
It is therefore a homomorphism and an isomorphism, making precise the claimed structural identity.

\exercise{11.52}{Detecting a noncyclic finite abelian group}
Suppose first $G$ contains $\Z_p^2$. This subgroup is not cyclic: every nonidentity element has order $p$, while the group has order $p^2$. Every subgroup of a cyclic group is cyclic, so $G$ cannot be cyclic.
Conversely, decompose finite abelian $G$ into its prime-power cyclic factors. If each prime occurred in at most one factor, these factor orders would be pairwise coprime, and their product would be cyclic. Since $G$ is not cyclic, some prime $p$ occurs in at least two factors, say $\Z_{p^a}$ and $\Z_{p^b}$. Their order-$p$ subgroups generated by $p^{a-1}$ and $p^{b-1}$ have direct product $\Z_p^2$, which embeds in $G$. This proves the equivalence.

\exercise{11.53}{Element orders in a prime-power group}
If $|G|=p^r$, then for any $g\in G$, Lagrange's theorem says $\operatorname{ord}(g)$ divides $p^r$. The positive divisors are exactly $p^s$ for $0\le s\le r$, so every element order is a power of $p$ (including $p^0=1$). This proof did not use commutativity; therefore the abelian hypothesis can be dropped.

\exercise{11.54}{Cancellation for finitely generated abelian groups}
By classification, a finitely generated abelian group is specified up to isomorphism by its free rank and, for every prime $p$ and exponent $j\ge1$, the multiplicity of the factor $\Z_{p^j}$. In a direct product these ranks and multiplicities add. The isomorphism $G\times K\cong H\times K$ therefore gives
\[r_G+r_K=r_H+r_K,\qquad m_G(p,j)+m_K(p,j)=m_H(p,j)+m_K(p,j).
\]
All quantities are finite nonnegative integers, since the groups are finitely generated. Cancelling the common terms yields equal free ranks and all equal prime-power multiplicities for $G,H$. The uniqueness part of the theorem then gives $G\cong H$.
